\begin{solapartat}
\punts{0,2} Segons la llei de gravitació universal, el mòdul de la força sobre un objecte de massa $m$ a la superfície d'un objecte astronòmic esfèric de radi $R$ i massa $M$ s'expressa com:
\[ F = G\,\frac{mM}{R^2} \]

\punts{0,1} I el mòdul de la força que experimenta un objecte de massa $m$ sota l'acció d'un camp gravitatori d'intensitat $g$ és
\[ F = mg \]
Igualant les dues expressions i dividint als dos costats per $m$ obtenim:
\[ g = G\,\frac{M}{R^2} \]

\punts{0,2} El radi de Didymos és $R_D = \frac{781}{2} = 390,5\ \mathrm{m}$ i la seva massa és:
\[ M_D = V\cdot\rho = \frac{4}{3}\pi R_D^3\rho = \frac{4}{3}\pi\cdot 390,5^3\cdot 2146 = 5,35\times 10^{11}\ \mathrm{kg} \]

\punts{0,35} Per tant obtenim:
% DUBTE: la pauta dona les unitats de g com m^2/s (haurien de ser m/s^2); es transcriu tal com és.
\[ g = G\,\frac{M_D}{R_D^2} = 6,67\times 10^{-11}\,\frac{5,35\times 10^{11}}{390,5^2} = 2,34\times 10^{-4}\ \frac{\mathrm{m^2}}{\mathrm{s}} \]

\punts{0,4} Segons la llei de gravitació universal, el mòdul de la força és
\[ F = G\,\frac{mM_D}{R^2} = 6,67\times 10^{-11}\,\frac{4,42\times 10^{10}\cdot 5,35\times 10^{11}}{1120^2} = 1,26\times 10^{6}\ \mathrm{N} \]
\end{solapartat}

\begin{solapartat}
\punts{0,2} Segons la llei de gravitació universal, el mòdul de la força sobre Dimorphos, massa $m$, degut a la atracció de Didymos és:
\[ F = G\,\frac{mM_D}{R^2} \]

\punts{0,2} I la segona llei de Newton estableix que: $\vec{F} = m\vec{a}$

\punts{0,2} D'altra banda, considerant que el satèl·lit descriu un moviment circular uniforme al voltant de Didymos, la seva acceleració centrípeta és: $a = \frac{v^2}{R}$

\punts{0,25} Com sobre Dimorphos només actua la força de la gravetat:
\[ G\,\frac{mM_D}{R^2} = m\,\frac{v^2}{R} \;\Rightarrow\; v = \sqrt{G\,\frac{M_D}{R}} \]

\punts{0,4} Podem comprovar que la velocitat és inversament proporcional a l'arrel quadrada del radi de l'òrbita, per tant, si aquest radi disminueix, el seu invers augmenta i, per tant, la velocitat també serà més gran, és a dir, quan el radi orbital disminueix la velocitat orbital augmenta.

També s'hauria pogut justificar a partir de la tercera llei de Kepler:
\[ \frac{R^3}{T^2} = G\,\frac{M_D}{4\pi^2} \ \text{i}\ T = \frac{2\pi R}{v} \;\Rightarrow\; v = \sqrt{G\,\frac{M_D}{R}} \]
\end{solapartat}
